\section{Fonctions}
\subsection{Parité d'une fonction}
Dans plusieurs contextes, connaître la parité d'une fonction peut être très utile afin de simplifier les calculs. On dit qu'une fonction $f(x)$ est paire si

\begin{equation}
    f(-x) = f(x)
\end{equation}

et on dit qu'elle est impaire dans le cas où

\begin{equation}
    f(-x) = -f(x)    
\end{equation}
  
Par exemple, la fonction $\cos(x)$ est une fonction paire et la fonction $\sin(x)$ est une fonction impaire. Sans grande surprise, le produit de deux fonctions paires $f(x)$ et $g(x)$ donne aussi une fonction paire. En effet, si on note $h(x) = f(x)g(x)$, alors

\begin{equation*}
    h(-x) = f(-x)g(-x) = f(x)g(x) = h(x)
\end{equation*}

Par ailleurs, si $f(x)$ et $g(x)$ sont deux fonctions impaires, leur produit donne encore une fois une fonction paire par une démonstration similaire. Toutefois, sans perte de généralité, si $f(x)$ est une fonction paire et $g(x)$ une fonction impaire, le produit donne une fonction impaire.

\begin{equation*}
    h(-x) = f(-x)g(-x) = f(x)(-1)g(x) = -h(x)
\end{equation*}

La parité d'une fonction peut nous éviter beaucoup de troubles lorsqu'on calcule son intégrale sur un intervalle symétrique. Pour une valeur $a$ du domaine et une fonction paire $f(x)$, on trouve que 

\begin{equation*}
    \int_{-a}^{a}f(x)dx = \int_{-a}^{0}f(x)dx + \int_{0}^{a}f(x)dx = \int_{-a}^{0}f(-x)dx + \int_{0}^{a}f(x)dx = -\int_{a}^{0}f(t)dt + \int_{0}^{a}f(x)dx
\end{equation*}
\begin{equation*}
    = \int_{0}^{a}f(t)dt + \int_{0}^{a}f(x)dx = 2\int_{0}^{a}f(x)dx
\end{equation*}

où on utilise la substitution $t = -x$. Bien qu'on intègre sur $t$ et sur $x$, les deux intégrales sont complètement équivalentes, car $t$ et $x$ sont des variables bidons. Ainsi, on se retrouve avec la même intégrale en double. Cependant, si $f(x)$ est impaire, on a 

\begin{equation*}
    \int_{-a}^{a}f(x)dx = \int_{-a}^{0}f(x)dx + \int_{0}^{a}f(x)dx = -\int_{-a}^{0}f(-x)dx + \int_{0}^{a}f(x)dx = \int_{a}^{0}f(t)dt + \int_{0}^{a}f(x)dx
\end{equation*}
\begin{equation*}
    = -\int_{0}^{a}f(t)dt + \int_{0}^{a}f(x)dx = 0
\end{equation*}

Encore une fois, on a utilisé la substitution $t = -x$ et la somme des intégrales s'annulent du fait qu'elles sont équivalentes. 

Finalement, on peut montrer que toute fonction $f(x)$ peut s'écrire comme la somme d'une fonction paire $g(x)$ et une fonction impaire $h(x)$.

\begin{equation*}
    f(x) = \frac{1}{2}f(x) + \frac{1}{2}f(x) + \frac{1}{2}f(-x) - \frac{1}{2}f(-x) = \left(\frac{1}{2}f(x) + \frac{1}{2}f(-x)\right) + \left(\frac{1}{2}f(x) - \frac{1}{2}f(-x)\right) = g(x) + h(x)
\end{equation*}

On vérifie que $g(x)$ est paire et que $h(x)$ est impaire.

\begin{equation*}
    g(-x) = \frac{1}{2}f(-x) + \frac{1}{2}f(-(-x)) = \frac{1}{2}f(-x) + \frac{1}{2}f(x) = g(x)
\end{equation*}
\begin{equation*}
    h(-x) = \frac{1}{2}f(-x) - \frac{1}{2}f(-(-x)) = \frac{1}{2}f(-x) - \frac{1}{2}f(x) = -\left(-\frac{1}{2}f(-x) + \frac{1}{2}f(x)\right) = -h(x)
\end{equation*}

Ainsi, $f(x) = g(x) + h(x)$ s'écrit comme la somme d'une fonction paire et d'une fonction impaire. Par ailleurs, cette décomposition est unique.

\subsection{Fonction périodique}
Une fonction $f(x)$ définie sur l'ensemble des réels est périodique et de période $T$ si

\begin{equation}
    f(x + T) = f(x)
\end{equation}

Autrement dit, une fonction périodique se répète après une période $T$. Par exemple, les fonctions $\cos(x)$ et $\sin(x)$ sont des fonctions périodiques ayant une période de $2\pi$. Un résultat important pour l'intégrale d'une fonction périodique est que 

\begin{equation}
    \int_{0}^{T}f(x)dx = \int_{a}^{a+T}f(x)dx
\end{equation}

En effet,

\begin{equation*}
    \derivative{}{a}\left(\int_{a}^{a+T}f(x)dx\right) = \derivative{}{a}\left(F(a+T) - F(a)\right) = \derivative{F(a+T)}{(a+T)} \derivative{(a+T)}{a} - \derivative{F(a)}{a} = f(a+T) - f(a) 
\end{equation*}
\begin{equation*}
    = f(a) - f(a) = 0
\end{equation*}

ce qui indique que l'intégrale qui vient d'être dérivée vaut la même chose peu importe la valeur de $a$. Donc, on en déduit (1.4) en prenant $a=0$.

\subsection{Fonctions orthogonales}
Pour deux vecteurs $\vec{u}$ et $\vec{v}$, ces derniers sont orthogonaux si leur produit scalaire est nul. 

\begin{equation*}
    \vec{u} \cdot \vec{v} = \sum_{i}^{}u_iv_i = 0
\end{equation*}

En quelque sorte, on peut dire qu'une fonction est un vecteur de taille infinie. Alors, le produit scalaire de fonctions, si on se fie à la définition pour les vecteurs, consisterait à sommer la valeur de leur produit à chaque point d'un domaine commun. Cependant, la somme ici est infinie et on discrétise le produit par un élément de largeur $dx$. Au final, on se retrouve à faire une intégrale. Ainsi, des fonctions $f(x)$ et $g(x)$ orthogonales respectent 

\begin{equation}
    \int_{a}^{b}f(x)g(x)dx = 0
\end{equation}

Par exemple, les fonctions $f(x) = 1$ et $ g(x) = x$ sur l'intervalle $[-1,1]$ sont des fonctions orthogonales puisque

\begin{equation*}
    \int_{-1}^{1}(1 \cdot x) dx = \int_{-1}^{1}x dx = 0 
\end{equation*}

par la parité de $x$. Si maintenant plusieurs fonctions $\{f_1, ..., f_n\}$ sont orthogonales les unes aux autres, on peut dire qu'il s'agit d'un ensemble orthogonal de fonctions respectant 

\begin{equation*}
    \int_{a}^{b}f_n(x)f_m(x)dx = 0 \ \ \ \forall n \neq m
\end{equation*}